% SAM 2 : une architecture volontairement asymétrique.
% Encodeur d'image lourd (calculé une fois) / encodeur d'invites et
% décodeur de masque légers (rejoués à chaque interaction).
                \begin{tikzpicture}[xscale=1.0, yscale=1.0]
                    \useasboundingbox (-0.35, -2.45) rectangle (12.85, 2.35);

                    \tikzset{lourd/.style={draw=gris_fonce_inria, line width=1.0pt, rounded corners=3pt,
                                           fill=gray!12, inner sep=4pt, align=center, font=\scriptsize,
                                           text width=2.55cm, minimum height=1.45cm}}
                    \tikzset{leger/.style={draw=inria-rouge, line width=1.0pt, rounded corners=3pt,
                                           fill=inria-rouge!7, inner sep=3pt, align=center, font=\scriptsize,
                                           text width=2.25cm, minimum height=0.90cm}}
                    \tikzset{donnee/.style={draw=inria-2024-bleu-canard, line width=0.8pt, rounded corners=2pt,
                                            fill=inria-2024-bleu-canard!12, inner sep=3pt, align=center,
                                            font=\scriptsize, text width=1.60cm, minimum height=0.80cm}}
                    \tikzset{fl/.style={-{Latex[length=2.2mm]}, line width=1.0pt, color=gray!55}}

                    % --- chaîne principale ---
                    \node[donnee] (img)  at (0.85, 1.10) {image};
                    \node[lourd]  (enc)  at (3.65, 1.10) {\textbf{encodeur d'image}\\{\tiny transformeur de vision\\(\emph{Hiera}), gelé}};
                    \node[donnee] (emb)  at (6.35, 1.10) {plongement\\$64\times64$};
                    \node[leger]  (dec)  at (9.20, 1.10) {\textbf{décodeur\\de masque}};
                    \node[donnee] (msk)  at (11.85, 1.10) {masque};

                    \draw[fl] (img) -- (enc);
                    \draw[fl] (enc) -- (emb);
                    \draw[fl] (emb) -- (dec);
                    \draw[fl] (dec) -- (msk);

                    % --- branche des invites ---
                    \node[donnee] (pts) at (3.65, -1.15) {points, boîte,\\masque grossier};
                    \node[leger]  (pe)  at (6.35, -1.15) {\textbf{encodeur\\d'invites}};
                    \draw[fl] (pts) -- (pe);
                    \draw[fl] (pe.east) -- (9.20,-1.15) -- (dec.south);

                    % --- annotations de coût ---
                    \node[font=\scriptsize, text=gris_fonce_inria, anchor=north, align=center] at (3.65, 0.20)
                        {calculé \textbf{une fois} par image};
                    \node[font=\scriptsize, text=inria-rouge, anchor=south, align=center] at (9.20, 1.95)
                        {rejoué à chaque invite (quelques dizaines de ms)};

                    \draw[inria-rouge, line width=0.7pt, densely dashed, rounded corners=3pt]
                        (2.35,-1.75) rectangle (7.65,-0.55);
                    \node[font=\scriptsize, text=inria-rouge, anchor=north, align=center] at (5.00, -1.85)
                        {c'est \textbf{ici} qu'une hiérarchie peut entrer sans rien casser};
                \end{tikzpicture}
